\section{Apparato sperimentale e obiettivi}
L'esperienza consiste nello studio della propagazione di onde elastiche trasversali su una corda vibrante tesa e nell'osservazione delle onde stazionarie che si instaurano quando la corda viene sollecitata da una forza periodica esterna. 

Quando una perturbazione si propaga lungo una corda tesa, le vibrazioni avvengono perpendicolarmente alla direzione di propagazione dell'onda, dando origine a onde trasversali, che possono essere descritte come funzioni di spazio $x$ e tempo $t$. Se la corda è fissata agli estremi, l'onda incidente viene riflessa e si propaga in direzione opposta. Se nel mentre sopraggiunge un'altra onda, esse si incontrano e lo spostamento totale del mezzo è dato dalla somma degli spostamenti dovuti alle singole onde, secondo il principio di sovrapposizione, e si ha:
\begin{equation}
    y(x,t) = 2A \sin\left(\frac{2\pi }{\lambda}x\right) \cos(2\pi f_n t)
    \label{eq:onda_stazionaria}
\end{equation}
dove $A$ è l'ampiezza massima, $\lambda$ è la lunghezza d'onda ed $f_n$ la sua frequenza. Si instaurano così onde stazionarie, per le quali la dipendenza spaziale è separata da quella temporale. Questo avviene solamente per particolari valori della frequenza della forzante. Sostituendo la lunghezza $L$ della corda nella relazione \eqref{eq:onda_stazionaria} e imponendo che gli estremi della corda siano nodi, dunque ad ampiezza nulla, segue che sulla corda sia presente un numero intero di semi-lunghezze d'onda.

L'apparato sperimentale è costituito da più corde, ciascuna con diversa massa lineare, da un sistema di sospensione, da un oscillatore elettromeccanico pilotato da un generatore di funzioni e da un sistema di carrucola e contrappesi per regolare la tensione, oltre a un metro a nastro e una bilancia analitica.
Gli obiettivi dell'esperienza sono i seguenti:

\begin{itemize}
    \item verificare sperimentalmente la formazione di onde stazionarie su una corda vibrante;
    \item studiare la dipendenza tra frequenza delle armoniche e numero armonico, tensione, lunghezza della corda e massa lineare;
    \item determinare la velocità di propagazione delle onde sulla corda e confrontarla con il valore previsto dal modello fisico.
\end{itemize}
\newpage